\documentclass[12pt]{article}

\usepackage[spanish]{babel}
\usepackage[utf8]{inputenc}
\usepackage[T1]{fontenc}
\usepackage{amsmath, amssymb}

\title{Evaluación 1 -- Parte II}
\author{Fundamentos de Matemática Básica}
\date{}

\begin{document}
	
	\maketitle
	
	\section*{Objetivo}
	\begin{itemize}
		\item Resolver paso a paso cada modelo algebraico.
		\item Simplificar expresiones, productos y fracciones racionales.
		\item Establecer restricciones cuando corresponda.
		\item Interpretar cada resultado en lenguaje técnico agrícola.
	\end{itemize}
	
	\section*{Problema 1}
	\subsection*{Modelos}
	\[
	M(x,y)=4xy-(x-2y)^2
	\]
	\[
	R(x,y)=(4x+y)(4y+x)
	\]
	\[
	F(x,y)=(x-2y)(x+2y)-(2x-y)^2
	\]
	
	\subsection*{Resultados}
	\[
	M(x,y)=-x^2+8xy-4y^2
	\]
	\[
	R(x,y)=4x^2+17xy+4y^2
	\]
	\[
	F(x,y)=-3x^2+4xy-5y^2
	\]
	
	\[
	T(x,y)=29xy-5y^2
	\]
	
	\section*{Problema 2}
	\[
	CT=(9a+7b)(a+4b)
	\]
	\[
	CM=(a+2b)(2a-3b)
	\]
	\[
	Cm=(2a-b)(3a+14b)
	\]
	
	\[
	CMO=a^2+17ab+48b^2
	\]
	
	\section*{Problema 3}
	\[
	R(x)=\frac{x^2-25}{x^2+5x}, \quad
	T(x)=\frac{25x^2+x}{x-5}
	\]
	
	\[
	M(x)=25x+1
	\]
	
	Restricciones:
	\[
	x \neq -5, \quad x \neq 0, \quad x \neq 5
	\]
	
	\section*{Problema 4}
	\[
	E=\frac{x^2}{4}
	\]
	
	Restricciones:
	\[
	x \neq -6, \quad x \neq 6, \quad x \neq 3, \quad x \neq -1
	\]
	
	\section*{Conclusión}
	La matemática permite modelar decisiones agrícolas como capacitación, análisis de suelo, eficiencia de riego y rendimiento productivo.
	
\end{document}